\documentclass[10pt,english,aspectratio=169]{beamer}
% Use notes or hide notes or show only notes or handout

\usetheme{default}

\usepackage{xstring}
\usepackage{pgfpages}
%\makeatletter
%\IfSubStr{\@classoptionslist}{handout}
%  {\pgfpagesuselayout{2 on 1}[letterpaper,border shrink=5mm]}
%  {}
%\makeatother

\usepackage{amsmath,amssymb,amsthm}
\usepackage{stmaryrd}
\usepackage{enumerate}
\usepackage{stfloats}
\usepackage{bbm}
\usepackage{pdfpages}
\usepackage{framed}

\usepackage[most]{tcolorbox}
\tcbset{highlight math style={enhanced,
  colframe=white,colback=yellow!15,arc=8pt,boxrule=1pt,
  }}
  
\usepackage{tikz,pgf,pgfplots}
\usepackage{algorithm,algorithmic}
\usepgflibrary{shapes}
\usetikzlibrary{%
  arrows,%
  arrows.meta,
  backgrounds,
  shapes.misc,% wg. rounded rectangle
  shapes.arrows,%
  shapes,%
  calc,%
  chains,%
  matrix,%
  positioning,% wg. " of "
  scopes,%
  decorations.pathmorphing,% /pgf/decoration/random steps | erste Graphik
  shadows,%
  backgrounds,%
  fit,%
  petri,%
  quotes
}

\tikzset{background rectangle/.style={
    fill=white,
  },
  use background/.style={    
    show background rectangle
  }
}

\setbeamersize{text margin left=10mm,text margin right=35mm}

\pgfplotsset{compat=1.12}

%\usetheme{Frankfurt}
%\usecolortheme{ldpc}
\useinnertheme{rounded}
\usecolortheme{whale}
\usecolortheme{orchid}

\newcommand{\ul}[1]{\underline{#1}}
\renewcommand{\Pr}{\mathbb{P}}

%% Setup slides and notes
\makeatletter
\IfSubStr{\@classoptionslist}{notes} { \IfSubStr{\@classoptionslist}{hide} {}{\IfSubStr{\@classoptionslist}{only} {}{\setbeameroption{show notes on second screen=right}}} }{}
\makeatother
%\setbeamertemplate{note page}{\pagecolor{yellow!5}\vfill\insertnote\vfill}

\newcommand{\getpdfpages}[2]{\begingroup
  \setbeamercolor{background canvas}{bg=}
  \addtocounter{framenumber}{1}
  \includepdf[pages={#1},%
  pagecommand={%
    \expandafter\def\expandafter\insertshorttitle\expandafter{%
      \insertshorttitle\hfill\insertframenumber\,/\,\inserttotalframenumber}}%
  ]{#2}
  \endgroup}

\newcommand{\backupbegin}{
   \newcounter{finalframe}
   \setcounter{finalframe}{\value{framenumber}}
}
\newcommand{\backupend}{
   \setcounter{framenumber}{\value{finalframe}}
}

 \setbeamercolor{bibliography entry author}{fg=black}
 \setbeamercolor{bibliography entry title}{fg=black}
 \setbeamercolor{bibliography entry location}{fg=black}
 \setbeamercolor{bibliography entry note}{fg=black}
 
 \setbeamerfont{bibliography item}{size=\footnotesize}
 \setbeamerfont{bibliography entry author}{size=\footnotesize}
 \setbeamerfont{bibliography entry title}{size=\footnotesize}
 \setbeamerfont{bibliography entry location}{size=\footnotesize}
 \setbeamerfont{bibliography entry note}{size=\footnotesize}
 \setbeamertemplate{bibliography item}{\insertbiblabel}
 
\newlength\tikzwidth
\newlength\tikzheight


\newcommand{\mc}[1]{\mathcal{#1}}
\newcommand{\mbb}[1]{\mathbb{#1}}
%\newcommand{\expt}{\mbb{E}}
%\newcommand{\dd}{\mathrm{d}}
\input{institution}

\newcommand{\Interior}[1]{\ensuremath{{#1}^{\circ}}}
\newcommand{\Closure}[1]{\ensuremath{\overline{#1}}}
\newcommand{\Complement}[1]{\ensuremath{{#1}^{c}}}

\newcommand{\Expect}{\ensuremath{\mathrm{E}}}
\newcommand{\vecnot}{\underline}
\newcommand{\RealNumbers}{\ensuremath{\mathbb{R}}}
\newcommand{\RationalNumbers}{\mathbb{Q}}
\newcommand{\ComplexNumbers}{\mathbb{C}}
\newcommand{\Real}{\mathrm{Re}}
\newcommand{\Span}{\mathrm{span}}
\newcommand{\Rank}{\mathrm{rank}}
\newcommand{\Nullity}{\mathrm{nullity}}
\newcommand{\Trace}{\mathrm{tr}}
\newcommand{\Diag}{\mathrm{diag}}
\newcommand{\dd}{\mathrm{d}}
\DeclareMathOperator*{\esssup}{ess\,sup}

% Use < , > inner product
\newcommand{\inner}[2]{{\left\langle #1 \mskip2mu , #2 \right\rangle}}
\newcommand{\tinner}[2]{{\langle #1 \mskip1mu , #2 \rangle}}

% Use < | > inner product
%\newcommand{\inner}[2]{{\left\langle #1 \mskip2mu \middle| \mskip2mu #2 \right\rangle}}
%\newcommand{\tinner}[2]{{\langle #1 \mskip1mu | \mskip1mu  #2 \rangle}}



\def\checkmark{\tikz\fill[scale=0.4](0,.35) -- (.25,0) -- (1,.7) -- (.25,.15) -- cycle;}
\def\greencheck{{\color{green}\checkmark}}
\def\scalecheck{\resizebox{\widthof{\checkmark}*\ratio{\widthof{x}}{\widthof{\normalsize x}}}{!}{\checkmark}}
\def\xmark{\tikz [x=1.4ex,y=1.4ex,line width=.2ex, red] \draw (0,0) -- (1,1) (0,1) -- (1,0);}
\def\redx{{\color{red}\xmark}}

\renewcommand{\footnotesep}{-2pt}


\begin{document}


\title{ECE 586: Vector Space Methods \\ Lecture 13 Flip Video: Normed Vector Spaces}
\author{Henry D. Pfister \\ \univ}
\date{}
%\date{August 20th, 2020}
%\maketitle

\setbeamertemplate{navigation symbols}{}

\begin{frame}[plain]
	\titlepage

	\note{
		\vspace{8mm}
		\begin{enumerate}
			\setlength\itemsep{3mm}
			\color{red}
			\item Welcome to the 13th video lecture for ECE 586, Vector Space Methods. \\[2mm]
			Today, we'll combine vector spaces with metric spaces by discussing normed vector spaces and Banach spaces.
		\end{enumerate}
	}
\end{frame}

\addtocounter{framenumber}{-1}
\setbeamertemplate{navigation symbols}{\textcolor{blue}{\footnotesize \insertframenumber ~/ \inserttotalframenumber}}




\begin{frame}<1-2>{3.5: Normed Vector Spaces}

Let $V$ be a vector space over a field $F$, where $F = \mathbb{R}$ or $\mathbb{C}$.

\begin{definition}<1->
A \textcolor{blue}{norm} on the vector space $V$ is a real-valued function $\left\| \cdot \right\| \colon V \rightarrow \RealNumbers$ that satisfies the following properties.
\begin{enumerate}
\item $\left\| \vecnot{v} \right\| \geq 0 \quad \forall \vecnot{v} \in V$;
equality holds if and only if $\vecnot{v} = \vecnot{0}$
\item $\left\| s \vecnot{v} \right\| = |s| \left\| \vecnot{v} \right\| \quad \forall \vecnot{v} \in V, s \in F$
\item $\left\| \vecnot{v} + \vecnot{w} \right\| \leq
\left\| \vecnot{v} \right\| + \left\| \vecnot{w} \right\| \quad \forall \vecnot{v}, \vecnot{w} \in V$.
\end{enumerate}
\end{definition}

\vspace{3mm}

\begin{uncoverenv}<2->
The concept of a norm is closely related to that of a metric.
For instance, a metric can be defined from any norm.\\[2mm]
For a norm $\left\| \cdot \right\|$ on the vector space $V$, the \textcolor{blue}{induced metric} is
\begin{equation*}
d \left( \vecnot{v}, \vecnot{w} \right)
= \left\| \vecnot{v} - \vecnot{w} \right\|.
\end{equation*}
\end{uncoverenv}

\note{
	\vspace{2mm}
	\begin{enumerate}[<alert@+>]
		\footnotesize
		\setlength\itemsep{3mm}
		\item A norm generalizes the idea of the length of a vector.  Read.  The third property is called the triangle inequality.
		\item Read.  Thus, every normed vector space is also a metric space and all of our results for metric spaces can be applied.
	\end{enumerate}
}

\end{frame}

\begin{frame}<1-2>{3.5: Examples of Normed Vector Spaces}

\begin{example}<1->[Standard Norms for Real/Complex Vector Spaces]
The following functions are examples of norms for $\mathbb{R}^n$ and $\mathbb{C}^n$:
\begin{enumerate}
\item the $l^1$ norm: $\left\| \vecnot{v} \right\|_1 = \sum_{i=1}^n |v_i|$
\item the $l^p$ norm: $\left\| \vecnot{v} \right\|_p = \left( \sum_{i=1}^n |v_i|^p \right)^{\frac{1}{p}}, \quad p \in (1,\infty)$
\item the $l^{\infty}$ norm: $\left\| \vecnot{v} \right\|_{\infty} = \max_{i\in \{1,\ldots, n\}} \{ |v_i| \}$
\end{enumerate}
\end{example}

\vspace{2mm}

\begin{example}<2->[Comparing Norms]
The same vector may have different lengths under different norms:
\[
\renewcommand{\arraystretch}{1.15}
\begin{array}{c|ccc}
\vecnot{v} & \|\vecnot{v}\|_1 & \|\vecnot{v}\|_2 & \|\vecnot{v}\|_\infty \\ \hline
(3,-4,12) & 19 & 13 & 12 \\
(1,\mathrm{i},1) & 3 & \sqrt{3} & 1 \\
(0,0,2) & 2 & 2 & 2
\end{array}
\]
\end{example}

\note{
	\vspace{2mm}
	\begin{enumerate}[<alert@+>]
		\footnotesize
		\setlength\itemsep{3mm}
		\item Read. The $l^2$ norm is the standard Euclidean length of a vector.
		\item Read. In the second row, $\mathrm{i}$ is the imaginary unit and its absolute value is one. The third row shows that these norms may also give the same length.
	\end{enumerate}
}

\end{frame}

\begin{frame}<1-3>{3.5: Norms for Function Spaces}

\begin{example}<1->[Standard Norms for Real/Complex Function Spaces]
For real or complex functions on $[a,b]$, replace sums by integrals:
\begin{enumerate}
\item the $L^1$ norm: $\left\| f(t) \right\|_1 = \int_a^b |f(t)| dt$
\item the $L^p$ norm: $\left\| f(t) \right\|_p = \left( \int_a^b |f(t)|^p dt \right)^{\frac{1}{p}}, \quad p \in (1,\infty)$
\item the $L^{\infty}$ norm: $\left\| f(t) \right\|_{\infty} = \esssup_{[a,b]} \{ | f(t) | \}$
\end{enumerate}
\end{example}

\begin{block}{Equality Almost Everywhere}<2->
\small
\begin{itemize}
\setlength\itemsep{0.5mm}
\item Integrals define \textcolor{blue}{seminorms} where $\|f\|_p=0$ iff $f=0$ \textcolor{blue}{almost everywhere (a.e.)}
\item Set of functions $=0$ a.e. form a \textcolor{blue}{closed subspace} in the seminorm topology
\item Equivalence relation $f\sim g$ when $\|f-g\|_p=0$ allows us to formally define $L^p$ as \textcolor{blue}{normed space of equivalence classes}.
\end{itemize}
\end{block}

\uncover<3->{In infinite dimensional spaces, some vectors/functions with finite values may have \textcolor{red}{infinite norm} and are \textcolor{blue}{excluded from the normed space} by convention.}

\note{
	\vspace{2mm}
	\begin{enumerate}[<alert@+>]
		\footnotesize
		\setlength\itemsep{3mm}
		\item Read. The essential supremum ignores changes on sets of zero length. For continuous functions on this interval, it equals the ordinary supremum.
		\item Read. For each $p$, we start with measurable functions for which the displayed expression is finite. A seminorm satisfies the norm axioms except that a nonzero function may have seminorm zero. Almost everywhere means except on a set of zero length. For example, changing a function at a single point does not change its equivalence class. The entire closed subspace of zero-a.e. functions becomes the zero vector, and the induced norm is well-defined on the equivalence classes. The standard construction uses the Lebesgue integral, but we will keep the usual function notation.
		\item Read. Otherwise, the norm would not be a real-valued function. An infinite-dimensional space does not mean that every vector has infinite norm; it means that no finite list spans the whole space.
	\end{enumerate}
}

\end{frame}

\begin{frame}<1-3>{3.5: Norms Versus Metrics}

\begin{example}
Consider vectors in $\RealNumbers^n$ with the Euclidean metric
\begin{equation*}
d \left( \vecnot{v}, \vecnot{w} \right)
= \sqrt{ (v_1 - w_1)^2 + \cdots + (v_n - w_n)^2 }.
\end{equation*}
Recall the bounded metric given by
\begin{equation*}
\bar{d} \left( \vecnot{v}, \vecnot{w} \right)
= \min \left\{ d \left( \vecnot{v}, \vecnot{w} \right), 1 \right\}.
\end{equation*}
\vspace{3mm}
Define $f \colon \RealNumbers^n \rightarrow \RealNumbers$ by
$f \left( \vecnot{v} \right) = \bar{d} \left( \vecnot{v}, \vecnot{0} \right)$.
Is the function $f$ a norm?

\begin{uncoverenv}<2->
By the properties of $\overline{d}$, we have
\begin{enumerate}
\item $\bar{d} \left( \vecnot{v}, \vecnot{0} \right) \geq 0 \quad \forall \vecnot{v} \in \mathbb{R}^n$; equality holds if and only if $\vecnot{v} = \vecnot{0}$ \addtocounter{enumi}{1}
\item $\bar{d} \left( \vecnot{v} + \vecnot{w}, \vecnot{0}\right) = \bar{d} \left( \vecnot{v}, -\vecnot{w} \right) \leq \bar{d} \left( \vecnot{v}, \vecnot{0} \right) + \bar{d} \left( \vecnot{w}, \vecnot{0} \right) \quad \forall \vecnot{v}, \vecnot{w} \in \mathbb{R}^n$.
\end{enumerate}
\end{uncoverenv}

\vspace{2mm}

\begin{uncoverenv}<3->
However, $\bar{d} \left( s \vecnot{v}, \vecnot{0} \right)$ is not always equal to $|s| \bar{d} \left( \vecnot{v}, \vecnot{0} \right)$.
For instance,
$\bar{d} \left( 2 \vecnot{e}_1, \vecnot{0} \right) = 1 < 2 \bar{d} \left( \vecnot{e}_1, \vecnot{0} \right)$.
Thus, the function $f \left( \vecnot{v} \right) = \bar{d} \left( \vecnot{v}, \vecnot{0} \right)$
is not a norm.
\end{uncoverenv}
\end{example}

\note{
	\vspace{2mm}
	\begin{enumerate}[<alert@+>]
		\footnotesize
		\setlength\itemsep{3mm}
		\item We have seen that every norm induces a metric.  But, does every metric define a norm?  Read.
		\item Read.  So, the function $f$ satisfies the first and third properties of a norm.
		\item Read.  In other words, not every metric is induced by a norm.
	\end{enumerate}
}

\end{frame}


\begin{frame}<1-3>{3.5: Properties of Normed Vector Spaces}

\begin{definition}<1->
A vector $\vecnot{v} \in V$ is called \textcolor{blue}{normalized} if $\left\| \vecnot{v} \right\| = 1$.
For any $\vecnot{v} \neq \vecnot{0}$, consider \vspace{-2mm}
\begin{equation*}
\vecnot{u} = \vecnot{v}\, / \left\| \vecnot{v} \right\|
\end{equation*}
with norm $\left\| \vecnot{u} \right\| = 1$.
A normalized vector is called a \textcolor{blue}{unit vector}.
\end{definition}

%\begin{block}
%We have not shown the $\ell^p$ and $L^p$ norms satisfy all the required properties.
%To prove the triangle inequality, one requires the Minkowski inequality which is stated in the notes.
%\end{block}

\begin{definition}<2->
For a vector space $V$ over $\mathbb{R}$ or $\mathbb{C}$, two norms (denoted $\|\cdot\|_{V}$ and $\|\cdot\|_{V'}$) are called \textcolor{blue}{equivalent} if there are positive reals $m,M$ such that, for all $\vecnot{v} \in V$,
\begin{align*}
m \| \vecnot{v} \|_{V'} & \leq \| \vecnot{v} \|_V \leq M \| \vecnot{v} \|_{V'}.
\end{align*}
\end{definition}

\begin{lemma}<3->
For finite-dimensional normed spaces, all norms are equivalent.
\end{lemma}

\note{
	\vspace{2mm}
	\begin{enumerate}[<alert@+>]
		\footnotesize
		\setlength\itemsep{3mm}
		\item Read.  The second property of a norm shows that dividing by the norm gives a vector whose norm equals one.
		\item Read.  Equivalent norms define the same open sets and hence the same notion of convergence.
		\item Read.  Thus, in finite dimensions, the choice of norm does not affect whether a sequence converges.  This is not true in infinite dimensions.
	\end{enumerate}
}

\end{frame}


\begin{frame}<1-4>[shrink=5]{3.5: Complete Normed Spaces}

%\begin{definition}<1->
%A vector $\vecnot{v} \in V$ is called \textcolor{blue}{normalized} if $\left\| \vecnot{v} \right\| = 1$.
%For any $\vecnot{v} \neq \vecnot{0}$, consider \vspace{-2mm}
%\begin{equation*}
%\vecnot{u} = \vecnot{v}\, / \left\| \vecnot{v} \right\|
%\end{equation*}
%with norm $\left\| \vecnot{u} \right\| = 1$.
%A normalized vector is called a \textcolor{blue}{unit vector}.
%\end{definition}

\begin{definition}<1->
A complete normed vector space is called a \textcolor{blue}{Banach space}.
\end{definition}

\begin{example}<2->
Vector space $\mathbb{R}^n$ (or $\mathbb{C}^n$) with any well-defined norm is a Banach space.
\end{example}

\begin{example}<3->
For $1 \leq p < \infty$, let $V = \ell^p$ be the Banach space of real sequences $(v_1,v_2,v_3,\ldots)$ with finite norm \vspace{-3mm}
\[ \|\vecnot{v} \|_p = \left(\sum_{i=1}^\infty |v_i |^p \right)^{1/p}. \]
\end{example}


\begin{example}<4->
The vector space of continuous functions $f \colon [a,b] \to \mathbb{R}$ is a Banach space under the supremum norm \vspace{-2.5mm}
\[ \left\| f(t) \right\| = \sup_{t\in [a,b]} |f(t)|. \vspace*{-0.5mm}\]
\end{example}

\note{
	\vspace{2mm}
	\begin{enumerate}[<alert@+>]
		\footnotesize
		\setlength\itemsep{3mm}
		\item Recall that a metric space is complete if every Cauchy sequence converges to a point in the space.  Read.  Completeness is defined using the induced metric.
		\item Read.  This follows from the completeness of the real numbers and the equivalence of all norms in finite dimensions.
		\item Read.  As we discussed, only sequences with finite norm are included.
		\item Read.  This works because the uniform limit of continuous functions is continuous.
	\end{enumerate}
}

\end{frame}

\begin{frame}<1-2>{3.5: Schauder Basis}

\begin{definition}<1->
A Banach space $V$ has a \textcolor{blue}{Schauder basis}, $\vecnot{v}_1, \vecnot{v}_2, \ldots$, if every $\vecnot{v} \in V$ can be written uniquely as
\[ \vecnot{v} = \sum_{i=1}^\infty s_i \vecnot{v}_i, \]
where convergence is determined by the norm topology.
\end{definition}

\begin{example}<1->
Let $V = \ell^p$, for $1 \leq p < \infty$, be the vector space of real sequences with finite $p$-norm.
The \textcolor{blue}{standard Schauder basis} is the countably infinite extension $\{\vecnot{e}_1 ,\vecnot{e}_2, \ldots \}$ of the standard basis.
\end{example}

\begin{block}{Caveat}<2->
A Banach space need not have a Schauder basis.
Thus, this cannot be assumed in general proofs for Banach spaces.
\end{block}

\note{
	\vspace{2mm}
	\begin{enumerate}[<alert@+>]
		\footnotesize
		\setlength\itemsep{3mm}
		\item Recall that a Hamel basis only allows finite linear combinations.  Now that we have a topology, we can also consider infinite linear combinations.  Read.
		\item Read.
	\end{enumerate}
}

\end{frame}

\begin{frame}<1-5>[shrink=5]{3.5: Convergence of Sums}


Banach space convergence via the induced metric $d(\vecnot{v}, \vecnot{w}) = \| \vecnot{v}-\vecnot{w} \|$.

\newcommand{\vtp}{\hspace{-0.02em}}
\begin{lemma}
If $\,\sum_{i=1}^\infty \| \vecnot{v}_i \| \vtp = \vtp a \vtp <\vtp \infty$, then $\vecnot{u}_n \vtp = \vtp \sum_{i=1}^n \vecnot{v}_i$ satisfies $\vecnot{u}_n \vtp \to \vtp \vecnot{u}$ with $\| \vecnot{u} \| \vtp \leq \vtp a$.
\end{lemma}
\vspace{-0.5mm}

\begin{proof}<2->
\begin{itemize}
\item<2-> Let $a_n = \sum_{i=1}^n \| \vecnot{v}_i \|$ and observe that, for $n\!>\!m$, \vspace{-2mm}
\[ |a_n - a_m| = \left| \sum_{i=1}^n \| \vecnot{v}_i \| - \sum_{i=1}^m \| \vecnot{v}_i \|  \right| = \sum_{i=m+1}^n \| \vecnot{v}_i \| \]
\vspace{-2mm}
\[ \|\vecnot{u}_n - \vecnot{u}_m \| = \left\| \sum_{i=1}^n \vecnot{v}_i - \sum_{i=1}^m \vecnot{v}_i  \right\| = \left\| \sum_{i=m+1}^n \vecnot{v}_i \right\| \leq \sum_{i=m+1}^n \| \vecnot{v}_i \| \]

\item<3-> Since $\sum_{i=1}^\infty \| \vecnot{v}_i \|$ converges in $\RealNumbers$, $a_n$ must be a Cauchy sequence \vspace{1mm}

\item<4-> Since $\|\vecnot{u}_n - \vecnot{u}_m\| \leq |a_n - a_m|$, $\vecnot{u}_n$ is also a Cauchy sequence \vspace{1mm}

\item<5-> Once $\vecnot{u}_n$ converges, the norm bound is given by the triangle inequality. \hfill \qedhere

\end{itemize}

\end{proof}

\note{
	\vspace{2mm}
	\begin{enumerate}[<alert@+>]
		\footnotesize
		\setlength\itemsep{3mm}
		\item Read.  This is the Banach space analog of the fact that absolutely convergent series of real numbers converge.
		\item Let $a_n$ equal the sum, from $i=1$ to $n$, of the norm of $\vecnot{v}_i$.  For $n>m$, the first equation shows that the absolute difference of the partial sums of the norms equals the sum of the intermediate terms.  The second equation shows that the norm of the difference of the vector partial sums is upper bounded by the sum of the norms of the intermediate terms.  The inequality follows from the triangle inequality.
		\item Read.
		\item Read.  Since a Banach space is complete, this Cauchy sequence converges to some vector $\vecnot{u}$.
		\item Read.  That is, the norm of $\vecnot{u}_n$ is at most $a$ for all $n$ and this bound also holds in the limit.
	\end{enumerate}
}

\end{frame}

\begin{frame}<1-3>[shrink=5]{3.5: Open and Closed Subspaces}

\begin{definition}<1->
A \textcolor{blue}{closed subspace} of a Banach space is a subspace that is a closed set in the topology generated by the norm.
\end{definition}
\vspace{-0.5mm}

\begin{theorem}<2->
All finite dimensional subspaces of a Banach space are closed.
\end{theorem}
\vspace{-0.5mm}

\begin{example}<3->[A span that is not closed]
Let $W = \{\vecnot{e}_1 ,\vecnot{e}_2 , \ldots \}$ be the standard Schauder basis of $\ell^p$, for $1 \leq p < \infty$.
The span of $W$ only includes finite linear combinations.
However, a sequence of finite linear combinations, like \vspace{-2mm}
\[ \vecnot{u}_n = \sum_{i=1}^n \frac{1}{i^2} \vecnot{e}_i , \vspace{-1.5mm} \]  converges to $\lim_{n\to \infty} \vecnot{u}_n$ with infinite support.
This limit is not a finite linear combination of the basis vectors. Thus, the span of $W$ is not closed.
\end{example}
\vspace{-0.5mm}

\visible<3->{Show convergence in live session.}

\note{
	\vspace{2mm}
	\begin{enumerate}[<alert@+>]
		\footnotesize
		\setlength\itemsep{3mm}
		\item Read.  Recall that a closed set contains all of its limit points.
		\item Read.  This is because every finite-dimensional normed space is complete.
		\item (read example, but not the final line)  Convergence will be shown in the live session.
	\end{enumerate}
}

\end{frame}

\begin{frame} \frametitle{Next Steps}

\begin{itemize}
\setlength\itemsep{5mm}
\item To continue studying after this video -- \vspace{2mm}

\begin{itemize}
 \setlength\itemsep{3mm}
 \item Try the required reading: Course Notes EF 3.5
 \item Or the recommended reading: LADR Ch. 6A
 \item Also, look at the problems in Assignment 5
\end{itemize}
\end{itemize}

\note{
	\vspace{8mm}
	\begin{enumerate}
		\setlength\itemsep{3mm}
		\color{red}
		\item Here are some options to continue learning this material. (read) \\ [2mm]  That's it for today.  So, I'll see you next time.
	\end{enumerate}
}

\end{frame}


\end{document}
